\section{Background}
\label{sec:background}
This section provides backgrounds on Transformer models and introduces parallelization schemes for LLM inference.

% \yx{Do we need a section for normal and AI-optimized FPGAs? We can introduce the AI Engine there.}

\subsection{Transformer Models}
\label{sec:transformer}
\begin{figure}[!htbp]
    \centering
    \includegraphics[width=0.8\linewidth]{figs/transformer-2stage.pdf}
    \caption{Transformer model. Red blocks represent linear operators, and blue blocks signify non-linear operators.}
    \label{fig:transformer}
\end{figure}

The Transformer model consists of both encoder and decoder blocks~\cite{vaswani2017transformer}.
Recent employment on LLMs mostly uses decoder-only models, which leverage an auto-regressive approach for text generation~\cite{radford2019gpt2,openai2023gpt4,touvron2023llama}.
We will mainly discuss decoder-only models in this paper, but since encoders and decoders share the core building blocks with subtle architectural variances, our approach can also be extended for encoder-only models~\cite{devlin2018bert,lan2020albert,liu2019roberta}.

As illustrated in Figure~\ref{fig:transformer}, generative inference of LLMs has two stages: prefill stage and decode stage~\cite{pope2022googleinf}.
In the prefill stage, the model takes in user prompts, normally a long sequence with $l_{\text{input}}$ tokens, goes through the whole Transformer model, and generates the first token.
In the decode stage, the model takes in the previously generated token and generates $l_{\text{gen}}$ new tokens one at a time in an auto-regressive way.
Since each token depends on the previously generated tokens, the decode stage is purely sequential.

We then go through the detailed model architecture.
The input tokens are first passed into an embedding layer that maps the discrete tokens into high-dimensional continuous representations while incorporating positional encoding for each token.
Subsequently, it generates a tensor (i.e., hidden states) of shape $(l,d)$, where $l$ represents sequence length, and $d$ is the size of hidden dimensions.
We omit the batch dimension to simplify the analysis, focusing solely on single-batch inference in this paper, but our approach can be easily extended to different batch sizes for LLM serving by adding an additional batch dimension~\cite{kwon2023vllm,li2023alpaserve}.

The hidden states then pass through a series of $N$ Transformer blocks.
Each Transformer block consists of two sublayers: a multi-head attention (MHA) module and a feed-forward network (FFN).
Residual connections and layer normalization (LayerNorm) functions are applied between these sublayers, although the specific order and application may vary across different models~\cite{xiong2020preln}.
The MHA module plays a crucial role in capturing token relationships within the input sequence.
The input is initially partitioned into $h$ segments, where $h$ corresponds to the number of attention heads.
To compute the attention scores for each head, the input sequence of length $l$ undergoes three linear projections: query, key, and value.
These projections, which are trainable, yield matrices $Q$, $K$, and $V$ respectively.
Attention scores are then computed using a scaled dot-product (SDP) operator between $Q$, $K$, and $V$, as specified by the formula:
\begin{equation}
\text{Attention}(Q,K,V)=\mathrm{softmax}\left(QK^T\Big/\sqrt{d_k}\right)V\,,
\end{equation}
where $d_k$ is the size of the hidden dimension.
The output from this operator involves $h$ outputs, which are subsequently concatenated and processed through an additional linear projection.
In the prefill stage, the generated $K$ and $V$ tensors will be stored as \emph{KV cache} and later be concatenated before SDP during the decode stage~\cite{pope2022googleinf}.

The FFN module comprises a linear layer followed by a non-linear activation function and another linear layer.
This module transforms the outputs of MHA into embedding matrices, which are then further processed by subsequent Transformer layers.

Finally, the output tensor will go through a softmax function to obtain a distribution.
The model will sample a token from this distribution and feed it into the decode stage.
For encoder-only models, there is only a prefill stage involved, and the distribution will be directly used for different downstream tasks like text classification~\cite{devlin2018bert,liu2019roberta,lan2020albert}.

In this paper, we only focus on analyzing the core Transformer blocks and accelerating them on FPGAs.
Embedding layers and output sampling~\cite{generation_strategies,chen2023speculative} require extensive random memory accesses, which may not be suitable for FPGA acceleration.
Also, they only take a small fraction of overall compute that does not affect the overall latency~\cite{kim2023survey}, so we leave them to execute on CPUs or GPUs as usual.

% \zz{we also need to properly discuss the distinction between the prefill and decode stages and their implications on hardware acceleration (we can put the implication part in a different section).}

\subsection{Parallelization Schemes}
\label{sub:parallelism_intro}

% \begin{figure}[!htbp]
    \centering
       
    \begin{subfigure}[b]{0.47\linewidth}
        \includegraphics[width=\linewidth]{figs/Data Parallelism.pdf}
        \caption{Data parallelism.}
    \end{subfigure}
    \begin{subfigure}[b]{0.49\linewidth}
        \includegraphics[width=\linewidth]{figs/Tensor Parallelism.pdf}
        \caption{Tensor parallelism}    
    \end{subfigure}
    
    \begin{subfigure}[t]{0.71\linewidth}
        \includegraphics[width=\linewidth]{figs/Pipeline Parallelism.pdf}    
        \caption{Pipeline parallelism}
    \end{subfigure}

    \caption{Parallelism schemes.}
    \label{fig:3d-parallelism}
\end{figure}

\begin{figure}[t]
    \centering
    \includegraphics[width=0.9\linewidth]{figs/3d-parallel.pdf}
    \caption{An example of tensor parallelism of a Transformer layer with two devices.
    TP rank is the unique identifier given to a device within a TP group.
    SM is the softmax function, LN is LayerNorm, and GL is the GeLU function.}
    % \yx{we need to explain ``TP rank'', or just use ``device'' instead.}
    \label{fig:3d-parallelism}
\end{figure}

% \zz{Parallelization Schemes}
% \zz{later we say we exploit/use/employ certain data/tensor/pipeline parallelism}
As model sizes continue to expand, it becomes increasingly common for a single device to be insufficient for accommodating the entire model.
Consequently, the exploration of diverse parallelization schemes within a device and across devices becomes necessary.
Parallelism techniques in deep learning can be roughly classified into data, tensor, and pipeline parallelism, together known as 3D parallelism~\cite{shoeybi2019megatron,narayanan2021megatronv2,korthikanti2022reducing,chen2024slapo}.
During the era of CNNs, data parallelism was the norm, involving the partitioning of input tensors along the batch dimension and their distribution across multiple devices~\cite{li2020ptddp,abadi2016tensorflow}.
DeepSpeed ZeRO~\cite{rajbhandari2020zero} and FSDP~\cite{fsdp} extended data parallelism by proposing a three-stage parallelism strategy that partitions optimizer states, gradients, and parameters to minimize memory usage.
However, this approach may incur high communication overheads during inference.

A more recent parallelism scheme for LLM inference is tensor parallelism (TP)~\cite{li2023alpaserve,aminabadi2022dsinf,shoeybi2019megatron,narayanan2021megatronv2,pope2022googleinf}, which distributes model parameters across multiple devices and conducts explicit collective operations to ensure model correctness.
Megatron-LM~\cite{shoeybi2019megatron} is the first to explore tensor parallelism for Transformer-based models, proving to be efficient in both training and inference due to relatively low communication costs.
As shown in Figure~\ref{fig:3d-parallelism}, tensor parallelism requires two \texttt{all\_reduce} operations inside a Transformer layer to ensure the results are correct.
Our accelerator design also explores tensor parallelism, as detailed in \S\ref{subsub:parallelism_schemes}.

Lastly, pipeline parallelism~\cite{narayanan2019pipedream,yang2021pipemare,narayanan2021pipe2bw} divides the model across network layers.
Multiple layers are grouped into a pipeline stage, and different stages are assigned to different devices.
Pipeline parallelism is typically employed across multiple nodes.
Since both tensor parallelism and pipeline parallelism handle only portions of the network model, they are collectively referred to as model parallelism.
We revisit these parallelization schemes in \S\ref{subsub:parallelism_schemes}.
